\documentclass[../../main.tex]{subfiles}

\begin{document}

\section{凸函数}\label{section:ch06-凸函数}

本节始终假设 $(a,b)$ 是有界或无界的开区间.

\begin{definition}[凸函数]\label{definition:ch06-凸函数}
$(a,b)$ 上的实值函数 $\varphi$ 称为\textbf{凸的}, 若对 $(a,b)$ 中的点 $x_{1},x_{2}$ 和每个满足 $0\leqslant\lambda\leqslant 1$ 的 $\lambda$,
\begin{align}
\varphi(\lambda x_{1}+(1-\lambda)x_{2})\leqslant\lambda\varphi(x_{1})+(1-\lambda)\varphi(x_{2}).\label{align:ch06-凸不等式}
\end{align}
\end{definition}

若我们看 $\varphi$ 的图像, 凸性不等式的几何阐述为 $(x_{1},\varphi(x_{1}))$ 与 $(x_{2},\varphi(x_{2}))$ 之间的弦上的每个点在 $\varphi$ 的图像的上方.

观察到对于 $(a,b)$ 中的两个点 $x_{1}<x_{2}$, $(x_{1},x_{2})$ 中的每个点 $x$ 可被表示为
\begin{align*}
x=\lambda x_{1}+(1-\lambda)x_{2},\quad\text{其中 }\lambda=\frac{x_{2}-x}{x_{2}-x_{1}}.
\end{align*}
因此凸性不等式可以写为对 $(a,b)$ 中的 $x_{1}<x<x_{2}$,
\begin{align*}
\varphi(x)\leqslant\left[\frac{x_{2}-x}{x_{2}-x_{1}}\right]\varphi(x_{1})+\left[\frac{x-x_{1}}{x_{2}-x_{1}}\right]\varphi(x_{2}).
\end{align*}
重新整理这些项, 该不等式也可改写为对 $(a,b)$ 中的 $x_{1}<x<x_{2}$,
\begin{align}
\frac{\varphi(x)-\varphi(x_{1})}{x-x_{1}}\leqslant\frac{\varphi(x_{2})-\varphi(x)}{x_{2}-x}.\label{align:ch06-弦斜率}
\end{align}
因此凸性也可几何地阐述为从 $(x_{1},\varphi(x_{1}))$ 到 $(x,\varphi(x))$ 的弦的斜率不大于从 $(x,\varphi(x))$ 到 $(x_{2},\varphi(x_{2}))$ 的弦的斜率.

\begin{proposition}\label{proposition:ch06-导数递增凸}
若 $\varphi$ 在 $(a,b)$ 上是可微的, 且它的导数 $\varphi'$ 是递增的, 则 $\varphi$ 是凸的. 特别地, $\varphi$ 是凸的, 若它在 $(a,b)$ 上有非负的二阶导数 $\varphi''$.
\end{proposition}

\begin{proof}
令 $x_{1},x_{2}$ 为 $(a,b)$ 中满足 $x_{1}<x_{2}$ 的点, 且令 $x$ 属于 $(x_{1},x_{2})$. 我们必须证明
\begin{align*}
\frac{\varphi(x)-\varphi(x_{1})}{x-x_{1}}\leqslant\frac{\varphi(x_{2})-\varphi(x)}{x_{2}-x}.
\end{align*}
为此, 应用中值定理于 $\varphi$ 在每个区间 $[x_{1},x]$ 和 $[x,x_{2}]$ 的限制以选取点 $c_{1}\in[x_{1},x]$ 和 $c_{2}\in[x,x_{2}]$, 使得
\begin{align*}
\varphi'(c_{1})=\frac{\varphi(x)-\varphi(x_{1})}{x-x_{1}}\quad\text{与}\quad\varphi'(c_{2})=\frac{\varphi(x_{2})-\varphi(x)}{x_{2}-x}.
\end{align*}
由于 $\varphi'$ 是递增的,
\begin{align*}
\frac{\varphi(x)-\varphi(x_{1})}{x-x_{1}}=\varphi'(c_{1})\leqslant\varphi'(c_{2})=\frac{\varphi(x_{2})-\varphi(x)}{x_{2}-x}.
\end{align*}
\end{proof}

\begin{example}\label{example:ch06-凸函数例子}
以下三个函数是凸的, 这是由于它们有非负的二阶导数:
\begin{align*}
\text{在 }(0,\infty)\text{ 上},\quad\text{对 }p\geqslant 1,\ \varphi(x)=x^{p};\\
\text{在 }(-\infty,\infty)\text{ 上},\quad\varphi(x)=e^{\alpha x};\\
\text{在 }(0,\infty)\text{ 上},\quad\varphi(x)=\ln(1/x).
\end{align*}
凸性的以下几何重述在证明凸函数的可微性质中有用.
\end{example}

\begin{lemma}[弦的斜率引理]\label{lemma:ch06-弦斜率}
令 $\varphi$ 在 $(a,b)$ 上是凸的. 若 $x_{1}<x<x_{2}$ 属于 $(a,b)$, 则对 $p_{1}=(x_{1},\varphi(x_{1}))$, $p=(x,\varphi(x))$, $p_{2}=(x_{2},\varphi(x_{2}))$, $\overline{p_{1}p}$ 的斜率 $\leqslant\overline{p_{1}p_{2}}$ 的斜率 $\leqslant\overline{p p_{2}}$ 的斜率.
\end{lemma}

\begin{proof}
重新整理不等式\eqref{align:ch06-弦斜率}的项, 将它写为以下两个等价形式:
\begin{align*}
\text{对 }(a,b)\text{ 中的 }x_{1}<x<x_{2},\quad\frac{\varphi(x_{1})-\varphi(x)}{x_{1}-x}\leqslant\frac{\varphi(x_{2})-\varphi(x_{1})}{x_{2}-x_{1}};\\
\text{对 }(a,b)\text{ 中的 }x_{1}<x<x_{2},\quad\frac{\varphi(x_{2})-\varphi(x_{1})}{x_{2}-x_{1}}\leqslant\frac{\varphi(x_{2})-\varphi(x)}{x_{2}-x}.
\end{align*}
\end{proof}

对于开区间 $(a,b)$ 上的函数 $g$ 和点 $x_{0}\in(a,b)$, 若
\begin{align*}
\lim_{h\to 0,h<0}\frac{g(x_{0}+h)-g(x_{0})}{h}\quad\text{存在且有限},
\end{align*}
我们将此极限记为 $g'(x_{0}^{-})$ 且称它为 $g$ 在 $x_{0}$ 的左导数. 类似地, 我们定义 $g'(x_{0}^{+})$, 且称它为 $g$ 在 $x_{0}$ 的右导数. 当然, $g$ 在 $x_{0}$ 是可微的当且仅当它在 $x_{0}$ 有相等的左导数与右导数. 凸函数的连续性与可微性从以下引理得到, 其证明从弦的斜率引理直接得到.

\begin{lemma}\label{lemma:ch06-单边导数}
令 $\varphi$ 为 $(a,b)$ 上的凸函数. 则 $\varphi$ 在每个点 $x\in(a,b)$ 有左导数和右导数. 此外, 对 $(a,b)$ 中满足 $u<v$ 的点 $u$, $v$, 这些单边导数满足以下等式:
\begin{align}
\varphi'(u^{-})\leqslant\varphi'(u^{+})\leqslant\frac{\varphi(v)-\varphi(u)}{v-u}\leqslant\varphi'(v^{-})\leqslant\varphi'(v^{+}).\label{align:ch06-单边导数}
\end{align}
\end{lemma}

\begin{corollary}\label{corollary:ch06-凸Lipschitz}
令 $\varphi$ 为 $(a,b)$ 上的凸函数. 则 $\varphi$ 是 Lipschitz 的, 因此 $\varphi$ 在 $(a,b)$ 的每个闭有界区间 $[c,d]$ 上是绝对连续的.
\end{corollary}

\begin{proof}
根据前一个引理, 对 $c\leqslant u<v\leqslant d$,
\begin{align}
\varphi'(c^{+})\leqslant\varphi'(u^{+})\leqslant\frac{\varphi(v)-\varphi(u)}{v-u}\leqslant\varphi'(v^{-})\leqslant\varphi'(d^{-}).\label{align:ch06-Lipschitz估计}
\end{align}
因此对所有 $u$, $v\in[c,d]$, $|\varphi(u)-\varphi(v)|\leqslant M|u-v|$, 其中 $M=\max\{|\varphi'(c^{+})|,|\varphi'(d^{-})|\}$. 因此 $\varphi$ 在 $[u,v]$ 上的限制是 Lipschitz 的. 闭有界区间上的 Lipschitz 函数是绝对连续的.
\end{proof}

我们从上面的推论和推论\cref{corollary:ch06-BV可微}推出, 定义在开区间上的任何凸函数在它的定义域上是几乎处处可微的. 事实上, 可以说更多.

\begin{theorem}\label{theorem:ch06-凸可微}
令 $\varphi$ 为 $(a,b)$ 上的凸函数. 则除可数个点外, $\varphi$ 是可微的且它的导数 $\varphi'$ 是递增函数.
\end{theorem}

\begin{proof}
我们从不等式\eqref{align:ch06-单边导数}推出, 函数
\begin{align*}
x\mapsto\varphi'(x^{-})\quad\text{与}\quad x\mapsto\varphi'(x^{+})
\end{align*}
是 $(a,b)$ 上的递增实值函数. 但根据定理\cref{theorem:ch06-单调连续}, 递增的实值函数除可数个点外是连续的. 因此, 除了 $(a,b)$ 的可数子集 $C$ 外, $\varphi$ 的左导数与右导数都是连续的. 令 $x_{0}$ 属于 $(a,b)\sim C$. 选取大于 $x_{0}$ 收敛于 $x_{0}$ 的点序列 $\{x_{n}\}$. 应用引理\cref{lemma:ch06-单边导数}, 其中 $x_{0}=u$ 而 $x_{n}=v$, 取极限得到
\begin{align*}
\varphi'(x_{0}^{-})\leqslant\varphi'(x_{0}^{+})\leqslant\varphi'(x_{0}^{-}).
\end{align*}
则 $\varphi(x_{0}^{-})=\varphi(x_{0}^{+})$, 因此 $\varphi$ 在 $x_{0}$ 是可微的. 为证明 $\varphi'$ 是 $(a,b)\sim C$ 上的递增函数, 令 $u$, $v$ 属于 $(a,b)\sim C$ 满足 $u<v$. 则根据引理\cref{lemma:ch06-单边导数},
\begin{align*}
\varphi'(u)\leqslant\frac{\varphi(v)-\varphi(u)}{v-u}\leqslant\varphi'(v).
\end{align*}
\end{proof}

令 $\varphi$ 为 $(a,b)$ 上的凸函数而 $x_{0}$ 属于 $(a,b)$. 对于实数 $m$, 过点 $(x_{0},\varphi(x_{0}))$ 的直线 $y=m(x-x_{0})+\varphi(x_{0})$, 称为 $\varphi$ 的图像在 $x_{0}$ 的\textbf{支撑线}, 若该直线总是落在 $\varphi$ 的图像的下方, 即若对所有 $x\in(a,b)$,
\begin{align*}
\varphi(x)\geqslant m(x-x_{0})+\varphi(x_{0}).
\end{align*}
从引理\cref{lemma:ch06-单边导数}得知这样的直线是支撑的当且仅当它的斜率 $m$ 落在 $\varphi$ 在 $x_{0}$ 的左右导数之间. 因此, 特别地, 在每个点总是至少存在一条支撑线. 该概念使我们能对以下不等式给出简短的证明.

\begin{theorem}[Jensen 不等式]\label{theorem:ch06-Jensen}
令 $\varphi$ 为 $(-\infty,\infty)$ 上的凸函数, $f$ 是 $[0,1]$ 上的可积函数, 且 $\varphi\circ f$ 在 $[0,1]$ 上也可积. 则
\begin{align}
\varphi\left(\int_{0}^{1}f(x)\,dx\right)\leqslant\int_{0}^{1}(\varphi\circ f)(x)\,dx.\label{align:ch06-Jensen}
\end{align}
\end{theorem}

\begin{proof}
定义 $\alpha=\int_{0}^{1}f(x)\,dx$. 选取 $m$ 落在 $\varphi$ 在点 $\alpha$ 的左右导数之间. 则 $y=m(t-\alpha)+\varphi(\alpha)$ 是 $\varphi$ 的图像在 $(\alpha,\varphi(\alpha))$ 的支撑线的方程. 因此,
\begin{align*}
\text{对所有 }t\in\mathbf{R},\quad\varphi(t)\geqslant m(t-\alpha)+\varphi(\alpha).
\end{align*}
由于 $f$ 在 $[0,1]$ 上是可积的, 它在 $[0,1]$ 上是 a.e. 有限的, 因此在该不等式中用 $f(x)$ 代替 $t$, 我们有
\begin{align*}
\text{对几乎所有 }x\in[0,1],\quad\varphi(f(x))\geqslant m(f(x)-\alpha)+\varphi(\alpha).
\end{align*}
对该不等式积分, 用 Lebesgue 积分的单调性与 $f$ 和 $\varphi\circ f$ 在 $[a,b]$ 上都是可积的假设, 得到
\begin{align*}
\int_{0}^{1}\varphi(f(x))\,dx\geqslant\int_{0}^{1}[m(f(x)-\alpha)+\varphi(\alpha)]\,dx=m\left[\int_{0}^{1}f(x)\,dx-\alpha\right]+\varphi(\alpha)=\varphi(\alpha).
\end{align*}
\end{proof}

关于 Jensen 不等式的假设, $\varphi\circ f$ 在 $[0,1]$ 上的可积性的一些说明是妥当的. 我们已经证明了凸函数是连续的, 因此第 3 章的命题 7 告诉我们若 $\varphi$ 是凸的且 $f$ 是可积的, 复合函数 $\varphi\circ f$ 是可测的. 若 $\varphi\circ f$ 非负, 则不必假设 $\varphi\circ f$ 是可积的, 这是由于若右边的积分等于 $+\infty$, 等式\eqref{align:ch06-Jensen}平凡成立. 在 $\varphi\circ f$ 不是非负的情形, 若存在常数 $c_{1}$ 和 $c_{2}$ 使得
\begin{align}
\text{对所有 }x\in\mathbf{R},\quad|\varphi(x)|\leqslant c_{1}+c_{2}|x|,\label{align:ch06-增长条件}
\end{align}
则我们从积分的比较判别法推出若 $f$ 在 $[0,1]$ 上可积, 则 $\varphi\circ f$ 在 $[0,1]$ 上也是可积的. 在缺乏增长假设\eqref{align:ch06-增长条件}的情形, 函数 $\varphi\circ f$ 在 $[0,1]$ 上可以是不可积的 (见习题 71).

\begin{exercise}\label{exercise:ch06-61}
证明 $(a,b)$ 上的实值函数 $\varphi$ 是凸的当且仅当对 $(a,b)$ 中的点 $x_{1},\cdots,x_{n}$ 和使得 $\sum_{k=1}^{n}\lambda_{k}=1$ 的非负数 $\lambda_{1},\cdots,\lambda_{n}$,
\begin{align*}
\varphi\left(\sum_{k=1}^{n}\lambda_{k}x_{k}\right)\leqslant\sum_{k=1}^{n}\lambda_{k}\varphi(x_{k}).
\end{align*}
用此直接证明对于 $f$ 是简单函数的 Jensen 不等式.
\end{exercise}

\begin{exercise}\label{exercise:ch06-62}
证明 $(a,b)$ 上的连续函数是凸的当且仅当
\begin{align*}
\text{对所有 }x_{1},x_{2}\in(a,b),\quad\varphi\left(\frac{x_{1}+x_{2}}{2}\right)\leqslant\frac{\varphi(x_{1})+\varphi(x_{2})}{2}.
\end{align*}
\end{exercise}

\begin{exercise}\label{exercise:ch06-63}
在一般的区间 $I$ 上的函数称为是凸的, 若它在 $I$ 上连续且\eqref{align:ch06-凸不等式}对所有 $x_{1},x_{2}\in I$ 成立. 闭有界区间 $[a,b]$ 上的凸函数必定在 $[a,b]$ 上是 Lipschitz 的吗?
\end{exercise}

\begin{exercise}\label{exercise:ch06-64}
令 $\varphi$ 在 $(a,b)$ 中的每个点有二阶导数. 证明 $\varphi$ 是凸的当且仅当 $\varphi''$ 是非负的.
\end{exercise}

\begin{exercise}\label{exercise:ch06-65}
假定 $a\geqslant 0$ 与 $b>0$. 证明对 $1\leqslant p<\infty$, 函数 $\varphi(t)=(a+bt)^{p}$ 在 $[0,\infty)$ 上是凸的.
\end{exercise}

\begin{exercise}\label{exercise:ch06-66}
对什么函数 $\varphi$, Jensen 不等式总是等式?
\end{exercise}

\begin{exercise}\label{exercise:ch06-67}
叙述和证明一般的闭有界区间 $[a,b]$ 上的 Jensen 不等式.
\end{exercise}

\begin{exercise}\label{exercise:ch06-68}
令 $f$ 为 $[0,1]$ 上的可积函数. 证明
\begin{align*}
\exp\left[\int_{0}^{1}f(x)\,dx\right]\leqslant\int_{0}^{1}\exp(f(x))\,dx.
\end{align*}
\end{exercise}

\begin{exercise}\label{exercise:ch06-69}
令 $\{\alpha_{n}\}$ 为非负数序列, 其和为 $1$, 而 $\{\zeta_{n}\}$ 是正数序列. 证明
\begin{align*}
\prod_{n=1}^{\infty}\zeta_{n}^{\alpha_{n}}\leqslant\sum_{n=1}^{\infty}\alpha_{n}\zeta_{n}.
\end{align*}
\end{exercise}

\begin{exercise}\label{exercise:ch06-70}
令 $g$ 为 $[0,1]$ 上的正可测函数. 证明 $\log\left(\int_{0}^{1}g(x)\,dx\right)\geqslant\int_{0}^{1}\log(g(x))\,dx$, 只要两边都有定义.
\end{exercise}

\begin{exercise}\label{exercise:ch06-71}
(Nemytskii) 令 $\varphi$ 为 $\mathbf{R}$ 上的连续函数. 证明: 若存在常数使得\eqref{align:ch06-增长条件}成立, 则只要 $f$ 在 $[0,1]$ 上可积, $\varphi\circ f$ 就在 $[0,1]$ 上可积. 接着证明: 若只要 $f$ 在 $[0,1]$ 上可积, $\varphi\circ f$ 就在 $[0,1]$ 上可积, 则存在常数 $c_{1}$ 和 $c_{2}$ 使得\eqref{align:ch06-增长条件}成立.
\end{exercise}

\end{document}